\addcontentsline{toc}{chapter}{Conclusions}
\chapter*{Conclusions}

DaQuVa demonstrates that a compact domain-specific language can express common
data-quality workflows over tabular data without requiring users to write SQL
or general-purpose Python scripts. The current prototype covers connection
declaration, table scanning, filtering, duplicate detection, merging, export,
controlled persistence, and an aggregate \texttt{summarize} command that emits a
proof table of before/after row counts for an end-to-end pipeline.

The latest implementation work improves both expressiveness and usability.
Compound filters now support \texttt{AND} / \texttt{OR} logic with clear
precedence, string predicates cover prefix, substring, and suffix matching,
and \texttt{sort} plus \texttt{limit} make exploratory top-$n$ queries possible
inside the DSL. The \texttt{summarize} command lets a pipeline emit an
inspectable proof of its effect to both CSV and SQLite, closing the
input--output--evidence loop expected of a data-quality tool. The interactive
REPL further supports incremental development by allowing users to execute
statements one at a time and inspect intermediate tables.

The runtime has also become safer and more predictable through clearer column
errors, safer SQLite replacement writes, normalized CSV string values, and
better handling of incompatible comparisons. These changes move DaQuVa from a
syntax prototype toward an executable data-quality tool that can be demonstrated
end to end on realistic cleaning scenarios.
